%double quote command
\newcommand\dquote[1]{``#1''}
\newcommand\dquoteit[1]{``\emph{#1}''}

%single quote command
\newcommand\squote[1]{`#1'}

% Segno di percentuale fuori da \SI (siunitx non lo definisce da solo)
\providecommand{\percent}{\%}

% =======================================================================================
% Macro per la notazione del lavoro.
% =======================================================================================

% RFD, LHS, RHS
\newcommand{\rfd}{\ensuremath{\mathit{rfd}}}
\newcommand{\LHS}{\ensuremath{\mathrm{LHS}}}
\newcommand{\RHS}{\ensuremath{\mathrm{RHS}}}
\newcommand{\NullCols}{\ensuremath{\mathrm{NullCols}}}
\newcommand{\thr}{\ensuremath{\theta}}

% Le tre condizioni derivate dal bytecode
\newcommand{\Vone}{\ensuremath{[\mathrm{V1}]}}
\newcommand{\Vtwo}{\ensuremath{[\mathrm{V2}]}}
\newcommand{\Vthree}{\ensuremath{[\mathrm{V3}]}}

% Strategie citate nel testo: monospazio per restare coerenti con il codice
\newcommand{\str}[1]{\texttt{#1}}

% Termini tecnici ricorrenti
\newcommand{\renuver}{\textsc{Renuver}}
\newcommand{\jar}{\texttt{RENUVER\_EDBT.jar}}

% Tabelle: intestazione compatta
\newcommand{\hd}[1]{\textbf{#1}}
